\chapter{Память SDRAM на 64 Мбит}
\label{ch:sdram}

\section{Зачем ПЛИС внешняя память}

Триггеры и блочная память BSRAM очень быстрые, но их мало. Для
изображения, звука или большого массива данных нужна память большего объёма.
Именно для этого в корпус GW2AR-18 встроена микросхема SDRAM на
\textbf{64 Мбит = 8 Мбайт}.

\begin{figure}[H]
\centering
\begin{tikzpicture}
\begin{axis}[
  xbar, bar width=5mm, width=0.8\textwidth, height=5.8cm,
  xmode=log, log origin=infty, xmin=1e3, xmax=3e8,
  xlabel={Объём, бит (логарифмическая шкала)},
  symbolic y coords={Триггеры, SSRAM, BSRAM, SDRAM},
  ytick=data, y tick label style={font=\sffamily\small},
  nodes near coords, nodes near coords style={font=\sffamily\small, anchor=west},
  point meta=explicit symbolic,
  label style={font=\sffamily\small}, x tick label style={font=\small},
  axis lines*=left, grid=major, grid style={FpgaGray!20}, enlarge y limits=0.2,
]
  \addplot[fill=FpgaOrange!70, draw=FpgaOrange] coordinates {
    (15552,Триггеры) [$\approx$1,9 Кбайт] (41472,SSRAM) [$\approx$5 Кбайт]
    (847872,BSRAM) [$\approx$103 Кбайт] (67108864,SDRAM) [8 Мбайт]};
\end{axis}
\end{tikzpicture}
\caption{Сколько данных можно хранить в разных видах памяти Tang Nano 20K}
\end{figure}

Например, один кадр изображения $640\times480$ точек по 16 бит на точку
занимает
\[
  640 \cdot 480 \cdot 16 = 4\,915\,200 \text{ бит} \approx 600 \text{ Кбайт},
\]
что в шесть раз больше всей блочной памяти ПЛИС, но легко помещается в
SDRAM: туда влезает больше десятка таких кадров.

\begin{diybox}[Сколько данных помещается во внутренней памяти]
\textbf{Задание.} Прикиньте, влезет ли ваш массив данных в блочную память, и посмотрите, как синтезатор размещает маленькую память.

\textbf{В Logra.} Карточка «На основе примера»: «RAM: запись и чтение» и «ROM из восьми слов». Вкладка «Код» $\to$ «Синтезировать»; вкладка Gowin или список ячеек в отчёте.

\textbf{Ожидаемый результат.} Массив \code{reg [7:0] mem [0:15]} размещается синтезатором как память (а не сотня триггеров), а для кадра $640\times 480\times 16$ бит ($\approx 4{,}9$ Мбит) блочной памяти (около 0,85 Мбит) не хватает.

\textbf{Как проверить.} Составьте таблицу «массив $\to$ объём в битах $\to$ влезает в BSRAM или нет» и сверьте с рисунком объёмов памяти в начале главы.

\stars{1} повторите: объём ПЗУ на 8 слов по 8 бит. \stars{2} измените: сделайте массив 256$\times$8 и посмотрите, как изменился отчёт. \stars{3} проект: этап «Охранная панель» --- журнал событий; критерий приёмки --- посчитано, сколько событий по 8 бит поместится во внутреннюю память и сколько в SDRAM.
\end{diybox}

\section{Как устроена DRAM}

Ячейка динамической памяти (DRAM) состоит всего из \textbf{одного
транзистора и одного конденсатора}. Заряженный конденсатор хранит 1,
разряженный --- 0. Поэтому DRAM такая плотная и дешёвая. Но у неё есть два
неудобства:
\begin{itemize}
  \item конденсатор постепенно \textbf{разряжается}, поэтому каждую строку
        нужно периодически перезаписывать (\term{регенерация}, refresh) --- все
        строки за 64~мс;
  \item чтение разрушает заряд, поэтому строку сначала копируют в
        быстрый буфер (\term{открывают}), а по окончании работы записывают
        обратно (\term{закрывают}).
\end{itemize}

\begin{figure}[H]
\centering
\begin{circuitikz}[scale=0.9, transform shape]
  \draw (3,1.5) node[nmos] (T) {};
  \draw (0,0 |- T.gate) node[left, font=\sffamily\small] {строка (WL)} -- (T.gate);
  \draw (T.drain) -- (T.drain |- 0,3.4) node[above, font=\sffamily\small] {столбец (BL)};
  \draw (T.source) to[C, l_=$C$] ++(0,-1.6) node[ground]{};
  \node[font=\sffamily\small, align=left, anchor=west] at (4.6,1.2) {1 бит = 1 транзистор + 1 конденсатор\\[1mm]
  \textcolor{FpgaRed}{заряд утекает} $\Rightarrow$ нужна регенерация};
\end{circuitikz}
\caption{Ячейка DRAM (упрощённо)}
\end{figure}

\begin{diybox}[Счётчик регенерации]
\textbf{Задание.} Все 2048 строк нужно обновить за 64~мс. Сделайте модуль, который выдаёт запрос \code{refresh\_req} с нужным периодом.

\textbf{В Logra.} Вкладка «Код»: счётчик с константой $N = 100\,\text{МГц} \cdot 31{,}25\,\text{мкс} = 3125$ (на модели возьмите $N=25$). Диаграммы: \code{cnt}, \code{refresh\_req}.

\textbf{Ожидаемый результат.} Один запрос на 3125 тактов (на модели --- на 25); \code{refresh\_req} длится один такт.

\textbf{Как проверить.} Диаграммы: на периоде $N$ ровно один импульс; $2048\cdot 31{,}25$~мкс $=64$~мс.

\stars{1} повторите: период 31,25 мкс. \stars{2} измените: частота SDRAM 133~МГц --- пересчитайте $N$. \stars{3} проект: этап «Охранная панель» --- регенерация журнала; критерий приёмки --- запрос регенерации не пропускается, даже когда идёт запись в журнал.
\end{diybox}

\section{Организация памяти}

SDRAM в GW2AR-18 имеет 32-разрядную шину данных и организована как
$4 \text{ банка} \times 2048 \text{ строк} \times 256 \text{ столбцов}
\times 32 \text{ бита}$:
\[
  4 \cdot 2048 \cdot 256 \cdot 32 = 67\,108\,864 \text{ бит} = 64 \text{ Мбит}.
\]

\begin{figure}[H]
\centering
\begin{tikzpicture}[x=1cm, y=1cm]
  \foreach \b/\c in {3/FpgaBlue!15, 2/FpgaBlue!25, 1/FpgaBlue!35, 0/FpgaBlue!50} {
    \begin{scope}[shift={(\b*0.45,\b*0.45)}]
      \fill[\c, draw=FpgaBlue] (0,0) rectangle (5,3.2);
      \node[font=\sffamily\scriptsize, anchor=north east] at (5,3.2) {банк \b};
    \end{scope}
  }
  \foreach \y in {0.4,0.8,...,3.0} \draw[FpgaBlue!40] (0,\y) -- (5,\y);
  \foreach \x in {0.5,1.0,...,4.6} \draw[FpgaBlue!40] (\x,0) -- (\x,3.2);
  \fill[FpgaOrange!80] (0,1.6) rectangle (5,2.0);
  \fill[FpgaRed] (2.5,1.6) rectangle (3.0,2.0);
  \draw[decorate, decoration={brace, amplitude=5pt}, thick] (-0.2,0) -- (-0.2,3.2) node[midway, left=6pt, font=\sffamily\small, align=right] {2048 строк\\\scriptsize адрес A[10:0]};
  \draw[decorate, decoration={brace, amplitude=5pt, mirror}, thick] (0,-0.2) -- (5,-0.2) node[midway, below=6pt, font=\sffamily\small, align=center] {256 столбцов, адрес A[7:0]};
  \node[font=\sffamily\small, align=left, anchor=west] at (7.4,2.6) {\textcolor{FpgaOrange}{\rule{3mm}{3mm}} открытая строка};
  \node[font=\sffamily\small, align=left, anchor=west] at (7.4,1.9) {\textcolor{FpgaRed}{\rule{3mm}{3mm}} одно 32-битное слово};
  \node[font=\sffamily\small, align=left, anchor=west] at (7.4,1.2) {банк выбирается BA[1:0]};
\end{tikzpicture}
\caption{Организация SDRAM: банк, строка, столбец}
\end{figure}

\begin{diybox}[Разрезать адрес на банк, строку и столбец]
\textbf{Задание.} Слово адресуется 21 битом. Разрежьте адрес на банк (2 бита), строку (11 бит) и столбец (8 бит).

\textbf{В Logra.} Вкладка «Код»: модуль с \code{input [20:0] addr} и \code{output [1:0] ba, output [10:0] row, output [7:0] col}; строка \code{assign \{ba, row, col\} = addr;}.

\textbf{Ожидаемый результат.} Для \code{addr = 0x1FFFFF} получается \code{ba = 3}, \code{row = 0x7FF}, \code{col = 0xFF}; для \code{addr = 0x000100} --- \code{ba = 0}, \code{row = 0}, \code{col = 0x00} и \code{row = 1} (следующая строка).

\textbf{Как проверить.} Диаграммы: задайте несколько адресов и сверьте поля; на RTL нет ни одного вентиля --- только провода, это переименование битов.

\stars{1} повторите: разрез для \code{0x1FFFFF}. \stars{2} измените: поменяйте порядок полей на \code{\{row, ba, col\}} и объясните, почему чередование банков ускоряет обмен. \stars{3} проект: этап «Охранная панель» --- адрес журнала; критерий приёмки --- счётчик адреса записи доходит до $2^{21}-1$ и возвращается в 0.
\end{diybox}

\section{Команды SDRAM}

В отличие от BSRAM, в SDRAM нельзя просто подать адрес и получить данные.
Микросхема управляется \term{командами}, которые кодируются комбинацией
сигналов \code{CS\#}, \code{RAS\#}, \code{CAS\#}, \code{WE\#} (решётка
означает активный низкий уровень) по фронту тактового сигнала.

\begin{table}[H]
\centering
\caption{Основные команды SDRAM}
\begin{tabular}{l c c c c l}
\toprule
\textbf{Команда} & \textbf{CS\#} & \textbf{RAS\#} & \textbf{CAS\#} & \textbf{WE\#} & \textbf{Действие} \\
\midrule
NOP        & 0 & 1 & 1 & 1 & ничего не делать \\
ACTIVE     & 0 & 0 & 1 & 1 & открыть строку в банке \\
READ       & 0 & 1 & 0 & 1 & прочитать из открытой строки \\
WRITE      & 0 & 1 & 0 & 0 & записать в открытую строку \\
PRECHARGE  & 0 & 0 & 1 & 0 & закрыть строку \\
AUTO REFRESH & 0 & 0 & 0 & 1 & регенерация \\
LOAD MODE  & 0 & 0 & 0 & 0 & настройка (CAS latency, длина пакета) \\
\bottomrule
\end{tabular}
\end{table}

\begin{figure}[H]
\centering
\begin{tikztimingtable}[timing/slope=0.1, xscale=1.9, yscale=1.4, timing/font=\sffamily\small, timing/rowdist=3.3ex]
  CLK      & 9{1L 1H} \\
  Команда  & 2D{ACTIVE} 2D{NOP} 2D{READ} 2D{NOP} 2D{NOP} 2D{PRE} 6D{NOP} \\
  Адрес    & 2D{строка} 2U 2D{столб.} 12U \\
  DQ[31:0] & 10Z 2D{D0} 2D{D1} 4Z \\
\extracode
  \draw[FpgaOrange, <->, thick] (1,-6.6) -- (5,-6.6) node[midway, below, font=\sffamily\scriptsize] {$t_\text{RCD}$};
  \draw[FpgaPurple, <->, thick] (5,-6.6) -- (9,-6.6) node[midway, below, font=\sffamily\scriptsize] {CAS latency = 2};
\end{tikztimingtable}
\caption{Чтение из SDRAM: открыть строку (ACTIVE), подождать $t_\text{RCD}$,
выдать команду READ, через CAS latency тактов получить данные}
\label{fig:sdram_read}
\end{figure}

\begin{diybox}[Кодировщик команд SDRAM]
\textbf{Задание.} Сделайте модуль, который по номеру команды выдаёт четыре сигнала \code{cs\_n}, \code{ras\_n}, \code{cas\_n}, \code{we\_n} по таблице команд.

\textbf{В Logra.} Вкладка «Код»: \code{case (cmd)} с ветками NOP, ACTIVE, READ, WRITE, PRECHARGE, REFRESH, LOAD MODE. Диаграммы: перебор \code{cmd = 0\ldots6}.

\textbf{Ожидаемый результат.} NOP $=0111$, ACTIVE $=0011$, READ $=0101$, WRITE $=0100$, PRECHARGE $=0010$, AUTO REFRESH $=0001$, LOAD MODE $=0000$ (порядок \code{cs\_n, ras\_n, cas\_n, we\_n}).

\textbf{Как проверить.} Диаграммы: сверьте все семь комбинаций с таблицей; ветка \code{default} должна выдавать NOP.

\stars{1} повторите: ACTIVE и READ. \stars{2} найдите ошибку: уберите ветку \code{default} и объясните, откуда на RTL взялась защёлка. \stars{3} проект: этап «Охранная панель» --- команды журнала; критерий приёмки --- все семь команд проверены на диаграмме без пропусков.
\end{diybox}

\section{Контроллер SDRAM}

Помнить обо всех командах и задержках при каждом обращении к памяти неудобно.
Поэтому между логикой пользователя и SDRAM ставят \term{контроллер} ---
модуль (обычно конечный автомат, как в проекте 5), который:
\begin{itemize}
  \item после включения питания инициализирует микросхему;
  \item превращает простые запросы «прочитай слово по адресу X» в
        последовательность команд ACTIVE $\to$ READ $\to$ PRECHARGE;
  \item сам вовремя выполняет регенерацию.
\end{itemize}

\begin{figure}[H]
\centering
\begin{tikzpicture}[node distance=24mm]
  \node[gblock, minimum height=24mm, minimum width=24mm] (user) {Ваша логика\\\scriptsize (видео, звук\ldots)};
  \node[hblock, right=of user, minimum height=24mm, minimum width=26mm] (ctrl) {Контроллер\\SDRAM\\\scriptsize конечный автомат};
  \node[oblock, right=of ctrl, minimum height=24mm, minimum width=20mm] (mem) {SDRAM\\64 Мбит};
  \draw[arr] ([yshift=6mm]user.east) -- node[above, lbl] {addr, wr, rd} ([yshift=6mm]ctrl.west);
  \draw[arr] ([yshift=0mm]user.east) -- node[above, lbl] {din[31:0]} ([yshift=0mm]ctrl.west);
  \draw[arr] ([yshift=-6mm]ctrl.west) -- node[below, lbl] {dout, busy} ([yshift=-6mm]user.east);
  \draw[arr] ([yshift=6mm]ctrl.east) -- node[above, lbl] {команды} ([yshift=6mm]mem.west);
  \draw[arr] ([yshift=0mm]ctrl.east) -- node[above, lbl] {A, BA} ([yshift=0mm]mem.west);
  \draw[arr, <->] ([yshift=-6mm]ctrl.east) -- node[below, lbl] {DQ[31:0]} ([yshift=-6mm]mem.west);
  \node[lbl, below=2mm of user] {простые запросы};
  \node[lbl, below=2mm of mem] {сложный протокол};
\end{tikzpicture}
\caption{Контроллер скрывает сложность протокола SDRAM}
\end{figure}

\subsection{Порты встроенной памяти}

В проекте Gowin встроенная SDRAM подключается к модулю верхнего уровня через
порты со специальными (зарезервированными) именами. Номера выводов в
\code{.cst} для них указывать не нужно: соединение находится внутри корпуса.

\begin{vcode}
module top (
    input  wire        clk,
    // --- Встроенная SDRAM (имена изменять нельзя!) ---
    output wire        O_sdram_clk,
    output wire        O_sdram_cke,
    output wire        O_sdram_cs_n,    // CS#
    output wire        O_sdram_ras_n,   // RAS#
    output wire        O_sdram_cas_n,   // CAS#
    output wire        O_sdram_wen_n,   // WE#
    output wire [3:0]  O_sdram_dqm,     // маска байтов
    output wire [10:0] O_sdram_addr,    // адрес строки/столбца
    output wire [1:0]  O_sdram_ba,      // номер банка
    inout  wire [31:0] IO_sdram_dq      // данные
);
\end{vcode}

\begin{infobox}[Где взять контроллер]
\begin{itemize}
  \item В Gowin EDA есть готовый IP-блок: \textbf{Tools $\to$ IP Core Generator
        $\to$ Memory Interface $\to$ SDRAM Controller}. Мастер сам
        сгенерирует модуль с простым интерфейсом.
  \item Существуют открытые контроллеры, написанные для Tang Nano 20K
        энтузиастами; их используют, например, проекты ретро-приставок NESTang и
        SNESTang.
  \item Хорошее упражнение повышенной сложности --- написать свой контроллер,
        следуя рис.~\ref{fig:sdram_read} и документации на микросхему.
\end{itemize}
\end{infobox}

\begin{diybox}[Скелет контроллера: ACTIVE, READ, PRECHARGE]
\textbf{Задание.} Опишите автомат, который по запросу \code{rd} выдаёт последовательность из рис.~\ref{fig:sdram_read}: ACTIVE, ожидание $t_\text{RCD}=2$, READ, ожидание CAS latency $=2$, PRECHARGE.

\textbf{В Logra.} Вкладка «Код»: состояния IDLE, ACT, WAIT1, RD, WAIT2, PRE (по образцу автомата светофора из проекта 5). Выход --- номер команды из предыдущего задания. Диаграммы: импульс \code{rd}.

\textbf{Ожидаемый результат.} После \code{rd} команды идут в порядке ACTIVE, NOP, READ, NOP, NOP, PRECHARGE и затем снова IDLE.

\textbf{Как проверить.} Диаграммы: поставьте курсор на READ и убедитесь, что между ACTIVE и READ ровно два такта. RTL: регистр состояния и логика переходов.

\stars{1} повторите: последовательность команд из рисунка. \stars{2} измените: сделайте $t_\text{RCD}=3$ параметром \code{parameter}. \stars{3} проект: этап «Охранная панель» --- контроллер журнала; критерий приёмки --- запись одного слова в журнал занимает последовательность ACTIVE, WRITE, PRECHARGE без нарушения задержек.
\end{diybox}

\section{Пропускная способность}

При тактовой частоте SDRAM, например, 100~МГц и шине 32 бита пиковая
скорость передачи равна
\[
  B = 100 \cdot 10^6 \ \text{Гц} \cdot 32 \ \text{бит} = 3{,}2 \ \text{Гбит/с} = 400 \ \text{Мбайт/с}.
\]
Реальная скорость ниже из-за открытия и закрытия строк и регенерации.
Чем длиннее пакеты последовательных адресов, тем ближе скорость к пиковой.

\begin{figure}[H]
\centering
\begin{tikzpicture}
\begin{axis}[
  width=0.82\textwidth, height=5.6cm,
  xlabel={Длина пакета, слов}, ylabel={Эффективная скорость, Мбайт/с},
  xmin=1, xmax=256, xmode=log, log basis x=2, ymin=0, ymax=420,
  xtick={1,2,4,8,16,32,64,128,256}, xticklabels={1,2,4,8,16,32,64,128,256},
  label style={font=\sffamily\small}, tick label style={font=\sffamily\small},
  grid=major, grid style={FpgaGray!20}, axis lines*=left,
]
  \addplot[very thick, FpgaBlue, samples at={1,2,4,8,16,32,64,128,256}, mark=*] {400*x/(x+6)};
  \addplot[thick, dashed, FpgaRed, domain=1:256] {400};
  \node[font=\sffamily\small, text=FpgaRed, anchor=north west] at (axis cs:1.3,388) {пиковая 400 Мбайт/с};
\end{axis}
\end{tikzpicture}
\caption{Оценка скорости обмена при 100 МГц: на каждый пакет приходится около
6 «служебных» тактов (ACTIVE, $t_\text{RCD}$, CAS latency, PRECHARGE)}
\end{figure}

\begin{ideabox}
SDRAM даёт много памяти (8 Мбайт), но требует контроллера: строки нужно
открывать и закрывать, а память регулярно регенерировать. Самая высокая
скорость получается, когда данные читаются и пишутся длинными
последовательными пакетами.
\end{ideabox}

\begin{diybox}[Измерьте эффективную скорость]
\textbf{Задание.} Посчитайте эффективную скорость обмена для пакетов длиной 1, 8, 64 и 256 слов.

\textbf{В Logra.} Вкладка «Диаграммы» на модели из предыдущего задания: добавьте счётчик тактов, пока \code{busy = 1}, и сравните с формулой $B(L)=400\cdot L/(L+6)$.

\textbf{Ожидаемый результат.} При $L=1$ получается около 57~Мбайт/с, при 8 --- 229, при 64 --- 366, при 256 --- 391~Мбайт/с: чем длиннее пакет, тем ближе скорость к пиковым 400.

\textbf{Как проверить.} Подставьте значения в формулу и нанесите точки на график главы.

\stars{1} повторите: $L=64$. \stars{2} измените: возьмите 8 служебных тактов вместо 6 и пересчитайте. \stars{3} проект: этап «Охранная панель» --- ёмкость и скорость журнала; критерий приёмки --- по одной формуле оценено, за какое время запишется журнал на 100\,000 событий.
\end{diybox}

\section{Вопросы и задачи}
\label{sec:tasks-sdram}

Ответы собраны в приложении~\ref{app:answers}.

\begin{taskbox}[Задачи к главе]
\begin{enumerate}
  \item \stars{1}~Сколько 32-битных слов хранит SDRAM из этой главы? Сколько разрядов должен иметь адрес слова?
  \item \stars{1}~Чему равна пиковая скорость при частоте 133~МГц и 32-разрядной шине?
  \item \stars{2}~За 64~мс нужно обновить все 2048 строк. Каков период команд AUTO REFRESH? Сколько тактов это при 100~МГц?
  \item \stars{2}~Найдите ошибку. Контроллер выдаёт READ на следующем такте после ACTIVE, хотя $t_\text{RCD}=2$.
        Что произойдёт? Как исправить автомат?
  \item \stars{3}~Проект. Напишите модуль, который разрезает 21-разрядный адрес слова на банк, строку и столбец
        и кодирует команды SDRAM. Критерий приёмки: на диаграмме все семь команд и несколько адресов дают
        значения из таблиц главы.
  \item \stars{3}~Проект «Охранная панель», этап «журнал». Опишите, как журнал событий будет храниться во внешней
        памяти. Критерий приёмки: указана разрядность счётчика адреса, период регенерации и оценка времени заполнения журнала.
\end{enumerate}
\end{taskbox}
